\documentclass[]{article}
\usepackage{geometry,tikz,lipsum,amsmath,amssymb,soul,xcolor,cancel}
\usepackage[dvipsnames]{xcolor}
\tikzset{every picture/.style={line width=0.75pt}} %set default line width to 0.75pt        
\numberwithin{equation}{section}

\NewDocumentCommand{\caja}{ m O{gray!40} O{gray!10} m }{
	\begin{figure}[h]
		\centering
		\begin{tikzpicture}
			\node[anchor=text, text width=\columnwidth-1.2cm, draw, rounded corners, line width=1pt, fill=#3, inner sep=5mm] (big) {\\#4};
			\node[draw, rounded corners, line width=.5pt, fill=#2, anchor=west, xshift=5mm] (small) at (big.north west) {#1};
		\end{tikzpicture}
	\end{figure}
}

\newcommand{\deriv}[3]{\left(\frac{\partial #1}{\partial #2} \right)_{#3}}

\newcommand{\dderiv}[3]{\left(\dfrac{\partial #1}{\partial #2} \right)_{#3}}

\NewDocumentCommand{\indUpDown}{ O{\Lambda} O{\mu} O{\nu} }{
	{#1}^{#2}_{\ #3}
}
\NewDocumentCommand{\indDownUp}{ O{\Lambda} O{\mu} O{\nu} }{
	{#1}_{#2}^{\ #3}
}

\newcommand{\LambdaMuNu}{\indUpDown}
\newcommand{\0}{\textcolor{gray}{0}}
\newcommand{\ei}{\mathbf{e}_i}
\newcommand{\ej}{\mathbf{e}_j}
\newcommand{\lc}{\epsilon_{ijk}}
\newcommand{\gr}{\nabla}
\newcommand{\di}{\nabla\cdot}
\newcommand{\ro}{\nabla\times}
\newcommand{\la}{\nabla^2}
\newcommand{\rv}{\mathbf{r}}
\newcommand{\pv}{\mathbf{p}}
\newcommand{\xv}{\mathbf{x}}
\newcommand{\yv}{\mathbf{y}}
\newcommand{\zv}{\mathbf{z}}
\newcommand{\xhat}{\hat{\xv}}
\newcommand{\yhat}{\hat{\yv}}
\newcommand{\zhat}{\hat{\zv}}
\newcommand{\ihat}{\hat{\mathbf{i}}}
\newcommand{\jhat}{\hat{\mathbf{j}}}
\newcommand{\khat}{\hat{\mathbf{k}}}
\newcommand{\rhat}{\hat{\rv}}
\newcommand{\that}{\hat{\mathbf{\theta}}}
\newcommand{\phat}{\hat{\mathbf{\varphi}}}
\newcommand{\rhhat}{\hat{\mathbf{\rho}}}
\newcommand{\nhat}{\hat{\mathbf{n}}}
\newcommand{\Rhat}{\hat{\mathbf{R}}}
\newcommand{\eo}{\epsilon_{0}}
\newcommand{\uo}{\mu_{0}}
\NewDocumentCommand{\ke}{ O{1}}{
	\frac{#1}{4\pi\eo}
}
\NewDocumentCommand{\ku}{ O{}}{
	\frac{\uo #1}{4\pi}
}
%\newcommand{\ke}{\frac{1}{4\pi\eo}}
%\newcommand{\ku}{\frac{\uo}{4\pi}}
\newcommand{\rdif}{\rv - \rv'}
\newcommand{\pd}[2]{\dfrac{\partial #1}{\partial #2}}
\newcommand{\lpd}[2]{\partial #1/\partial #2}
\newcommand{\spd}[2]{\frac{\partial #1}{\partial #2}}
\newcommand{\td}[2]{\dfrac{d #1}{d #2}}
\newcommand{\ltd}[2]{d #1/d #2}
\newcommand{\std}[2]{\frac{d #1}{d #2}}
\newcommand{\n}[1]{\mathbf{#1}}
\newcommand{\eff}[1]{#1_\mathrm{eff}}
\newcommand{\Ueff}{\eff{U}}

\newcommand{\wfrac}[3][3pt]{%
	\frac{\wfracterm{depth}{\dp}{#1}{#2}}{\wfracterm{height}{\ht}{#1}{#3}}%
}
\newcommand{\wfracterm}[4]{%
	\sbox0{$\displaystyle#4$}%
	\vrule width 0pt #1 \dimexpr #20+#3\relax
	\usebox{0}%
}

\newcommand{\pr}[1]{\left\langle #1 \right\rangle }

\newcommand{\qed}{\hfill\(\blacksquare\)}
\newcommand{\resp}{
	\\
	
	\textbf{Respuesta:} }
\newcommand{\celsius}{^\circ\mathrm{C}}

\newcommand{\GeV}{\,\text{GeV}}

%opening
\title{Certamen 2 Mecánica Clásica Avanzada}
\author{Simon Fonseca}

\begin{document}

\maketitle
\section{Problema 1}
A pointlike nonrelativistic particle of mass $m$ scatters in the time-dependent spherically symmetric potential
\begin{equation}
	U(r,t) = \frac{U_0}{r^2 + \kappa t}, ~~~~~~~ r = |\n{r}| = \sqrt{x^2 + y^2 + z^2}, ~~~~~~~ U_0, \kappa > 0 \label{eq:1.1}
\end{equation}

\begin{enumerate}
%%%%%%%%%%%%%%%%%%%%%%%%%%%%%%%%%%%%%%%%%%%%%%%%%%%%%%%%%%%%%%%%%%%%%%%%%%%%%%%%%%
	\item[(a)] Identify all continuous symmetries of the action. Analyze if there are symmetries of type
	\begin{equation}
		r \rightarrow \alpha r, ~~~~~ t \rightarrow \beta t, ~~~~~~~ \alpha, \beta = \mathrm{const.}
	\end{equation}
	which do not change the action or the equations of motion.
	\resp  El lagrangiano será
	\begin{equation}
		L = \frac{1}{2}m \left(\dot{r}^2 + r^2 \dot{\theta}^2 + r^2\dot{\phi}^2 \sin^2 \theta \right) - \frac{U_0}{r^2 + \kappa t}
	\end{equation}
	
	La variable $\phi$ es cíclica, por lo que $M_z$ será conservado. Aun más, el sistema es esféricamente simétrico (el potencial es central), por lo que todo el vector de momento angular se conservará:
	\begin{equation}
		\n{M} = \mathrm{const.}
	\end{equation}
	
	Reemplazando $\rv \rightarrow \alpha \rv, t\rightarrow \beta t$ en las derivadas,
	\begin{equation}
		\pd{r}{t} \rightarrow \pd{\alpha r}{\beta t} = \frac{\alpha}{\beta} \pd{r}{t}
	\end{equation}
	Luego en el lagrangiano
	\begin{align*}
		L' &= \frac{1}{2}m \left(\frac{\alpha^2}{\beta^2} \dot{r}^2 + \frac{\alpha^2}{\beta^2} r^2 \dot{\theta}^2 + \frac{\alpha^2}{\beta^2} r^2\dot{\phi}^2 \sin^2 \theta \right) - \frac{U_0}{\alpha^2 r^2 + \kappa \beta t}\\
		L' &= \frac{\alpha^2}{\beta^2} \left[\frac{1}{2}m \left(\dot{r}^2 + r^2 \dot{\theta}^2 + r^2\dot{\phi}^2 \sin^2 \theta \right) \right] - \frac{U_0}{\alpha^2 r^2 + \kappa \beta t}
	\end{align*}
	Al considerar el efecto sobre la acción,
	\begin{align*}
		S = \int L \, dt ~\rightarrow~ S' &= \int L' \, dt' \\
		&= \int \left(\frac{\alpha^2}{\beta^2} \left[\frac{1}{2}m \left(\dot{r}^2 + r^2 \dot{\theta}^2 + r^2\dot{\phi}^2 \sin^2 \theta \right)\right] - \frac{U_0}{\alpha^2 r^2 + \kappa \beta t} \right) \cdot \beta dt\\
		&= \int \left(\frac{\alpha^2}{\beta} \left[\frac{1}{2}m \left(\dot{r}^2 + r^2 \dot{\theta}^2 + r^2\dot{\phi}^2 \sin^2 \theta \right)\right] - \frac{U_0}{\frac{\alpha^2}{\beta} r^2 + \kappa t} \right)  \, dt
	\end{align*}
	Para que podamos usar el teorema de Noether la acción deberá ser invariante ($S' = S$), por lo que el factor $\frac{\alpha^2}{\beta}$ deberá cumplir
	\begin{equation*}
		\frac{\alpha^2}{\beta} = 1 ~~ \Rightarrow ~~ \beta = \alpha^2
	\end{equation*}
	\begin{equation*}
		\Rightarrow \rv \rightarrow \alpha \rv, ~~~~~ t\rightarrow \alpha^2 t, ~~~~~ \alpha = \mathrm{const.}
	\end{equation*}
	Podemos comprobar la invarianza reemplazando directamente en la acción:
	\begin{align*}
		S' &= \int \left[\left(\pd{\alpha r}{\alpha^2 t}\right)^2 + \alpha^2 r^2 \left(\pd{\theta}{\alpha^2 t}\right)^2 + \alpha^2 r^2 \left(\pd{\phi}{\alpha^2 t}\right)^2 \sin^2 \theta - \frac{U_0}{\alpha^2 r^2 + \kappa \alpha^2 t} \right] \cdot \alpha^2 dt\\
		&= \int \left(\frac{1}{\alpha^2} \left[\frac{1}{2}m \left(\dot{r}^2 + r^2 \dot{\theta}^2 + r^2\dot{\phi}^2 \sin^2 \theta \right)\right] - \frac{1}{\alpha^2} \frac{U_0}{ r^2 + \kappa t} \right) \cdot \alpha^2 dt\\
		&= \int \frac{\alpha^2}{\alpha^2} \left(\frac{1}{2}m \left(\dot{r}^2 + r^2 \dot{\theta}^2 + r^2\dot{\phi}^2 \sin^2 \theta \right) - \frac{U_0}{ r^2 + \kappa t} \right) \, dt\\
		&= \int \left(\frac{1}{2}m \left(\dot{r}^2 + r^2 \dot{\theta}^2 + r^2\dot{\phi}^2 \sin^2 \theta \right) - \frac{U_0}{ r^2 + \kappa t} \right) \, dt\\
		&= \int L \, dt = S
	\end{align*} 
	\qed
	
%%%%%%%%%%%%%%%%%%%%%%%%%%%%%%%%%%%%%%%%%%%%%%%%%%%%%%%%%%%%%%%%%%%%%%%%%%%%%%%%%%
	\item[(b)] Construct the integrals of motion which may exist according to Noether's theorem. Demonstrate that these quantities are conserved using the equations of motion.
	\resp El momento angular vendrá dado por
	\begin{equation}
		\n{M} = m\rv\times\dot{\rv} = \mathrm{const.}
	\end{equation}
	cuyas componentes en coordenadas cartesianas (estoy copiando esto desde la resolución de la Tarea 3) son
	\begin{align*}
		M_x &= -m r^2 \left( \dot{\theta}\sin\phi + \dot{\phi}\sin\theta\cos\theta\cos\phi \right)\\
		M_y &= m r^2 \left( \dot{\theta}\cos\phi - \dot{\phi}\sin\theta\cos\theta\sin\phi \right)\\
		M_z &= m r^2 \dot{\phi} \sin^2\theta
	\end{align*}
	
	Para la simetría de escala, escribiendo las transformaciones en notación con $\varepsilon$
	\begin{equation*}
		\alpha = e^{\varepsilon(t)} \approx 1 + \varepsilon(t)
	\end{equation*}
	tal que $\alpha$ es generado por cambios infinitesimales $\varepsilon$, los cambios de coordenadas y tiempo será, entonces
	\begin{equation*}
		\rv' \approx \rv + \delta \rv \approx \rv + \varepsilon \, \rv, ~~~~~ t' \approx t + \delta t \approx t + (1 + \varepsilon)^2 t \approx t + 2\varepsilon \, t
	\end{equation*}
	\begin{equation*}
		\delta \rv \approx \varepsilon \rv, ~~~~~ \delta t \approx 2\varepsilon t
	\end{equation*}
	\begin{align*}
		\Rightarrow d \rv' &= d \left(e^{\varepsilon(t)} \rv\right)\\
		&= e^{\varepsilon} d\rv + \rv e^{\varepsilon} d \varepsilon\\
		&= e^{\varepsilon} \dot{\rv} dt + \rv e^{\varepsilon} \dot{\varepsilon} dt \\
		&\approx (1 + \varepsilon) \dot{\rv} dt + \rv (1 + \varepsilon) \dot{\varepsilon} dt \\
		&\approx \underbrace{\dot{\rv} dt}_{d\rv} + \underbrace{\varepsilon \dot{\rv} dt}_{\mathcal{O}(\varepsilon)} + \underbrace{\rv \dot{\varepsilon} dt}_{\mathcal{O}(\dot{\varepsilon})} + \underbrace{\rv \varepsilon \dot{\varepsilon} dt}_{\mathcal{O}(\varepsilon \dot{\varepsilon})}\\
		&\\
		\Rightarrow d t' &= d \left(e^{2\varepsilon(t)} t\right)\\
		&= e^{2\varepsilon} dt + t e^{2\varepsilon} \cdot 2 d \varepsilon\\
		&= e^{2\varepsilon} dt + 2 t e^{2\varepsilon} \dot{\varepsilon} dt \\
		&\approx (1 + \varepsilon)^2 dt + 2 t (1 + \varepsilon)^2 \dot{\varepsilon} dt \\
		&\approx (1 + 2\varepsilon + \varepsilon^2) dt + 2 t (1 + 2\varepsilon + \varepsilon^2) \dot{\varepsilon} dt \\
		&\approx dt + \underbrace{2\varepsilon dt}_{\mathcal{O}(\varepsilon)} + \underbrace{\varepsilon^2 dt}_{\mathcal{O}(\varepsilon^2)} + \underbrace{2 t \dot{\varepsilon} dt}_{\mathcal{O}(\dot{\varepsilon})} + \underbrace{4 t \varepsilon \dot{\varepsilon} dt}_{\mathcal{O}(\varepsilon \dot{\varepsilon})} + \underbrace{2 t \varepsilon^2 \dot{\varepsilon} dt}_{\mathcal{O}(\varepsilon^2 \dot{\varepsilon})}
	\end{align*}
	Todos los términos $\varepsilon \dot{\varepsilon}, ~ \varepsilon^2, ~ \varepsilon^2 \dot{\varepsilon}$ son insignificantes, y después de esto deberemos diferenciar por $\dot{\varepsilon}$, por lo que solo nos interesarán los términos $\mathcal{O}(\dot{\varepsilon})$
	\begin{equation*}
		d \rv' \approx \left(\dot{\rv} + \dot{\varepsilon} \rv\right) dt, ~~~~~ dt' \approx \left(1 + 2 \dot{\varepsilon} t\right) dt
	\end{equation*}
	Luego, la velocidad
	\begin{equation*}
		\dot{\rv}' \equiv \frac{d \rv'}{d t'} \approx \frac{\left(\dot{\rv} + \dot{\varepsilon} \rv \right) dt}{\left(1 + 2 \dot{\varepsilon} t\right) dt}
	\end{equation*}
	Aproximando la división con \texttt{Mathematica},
	\begin{equation*}
		\dot{\rv}' \approx \dot{\rv} + \delta \dot{\rv} \approx \dot{\rv} + \dot{\varepsilon}\left(\rv - 2 \dot{\rv} t\right)
	\end{equation*}
	\begin{equation*}
		\delta \dot{\rv} \approx \dot{\varepsilon}\left(\rv - 2 \dot{\rv} t\right)
	\end{equation*}
	Luego, aproximando el lagrangiano a primer orden,
	\begin{align*}
		L'(\rv',\dot{\rv}',t') &= L'(\rv + \delta \rv, \dot{\rv} + \delta \dot{\rv}, t + \delta t)\\
		&\approx L(\rv,\dot{\rv},t) + \pd{L}{\rv} \cdot \delta\rv + \pd{L}{\dot{\rv}} \cdot \delta\dot{\rv} + \pd{L}{t} \delta t\\
		&\approx L(\rv,\dot{\rv},t) + \pd{L}{\rv} \cdot \varepsilon \rv + \pd{L}{\dot{\rv}} \cdot \dot{\varepsilon}\left(\rv - 2 \dot{\rv} t\right) + \pd{L}{t} 2 \varepsilon t
	\end{align*}
	De nuevo, solo nos interesarán los términos $\mathcal{O}(\dot{\varepsilon})$, por lo que, considerando que $\lpd{L}{\dot{\rv}} = m \dot{\rv}$
	\begin{equation*}
		L' \approx L + m \dot{\rv} \cdot \dot{\varepsilon}\left(\rv - 2 \dot{\rv} t\right) 
	\end{equation*}
	Por lo que la acción
	\begin{align*}
		S' &= \int L' \, dt'\\
		&\approx \int \left[L + m \dot{\rv} \cdot \dot{\varepsilon}\left(\rv - 2 \dot{\rv} t\right)\right] \, \left[1 + 2 \dot{\varepsilon} t\right] dt\\
		&\approx \int L \, dt + \int m \dot{\rv} \cdot \dot{\varepsilon} \left(\rv - 2 \dot{\rv} t\right) \, dt + \int 2 \dot{\varepsilon} t L \, dt + \int 2 m \dot{\rv} \cdot \dot{\varepsilon}^2 t \left(\rv - 2 \dot{\rv} t\right) \, dt
	\end{align*}
	El último término es $\mathcal{O}(\dot{\varepsilon}^2)$, por lo que es insignificante. Tendremos entonces
	\begin{equation*}
		S' \approx \int \left[L + m \dot{\rv} \cdot \dot{\varepsilon} \left(\rv - 2 \dot{\rv} t\right) + 2 \dot{\varepsilon} t L\right] \, dt
	\end{equation*}
	Y la variación de $L$ será
	\begin{equation*}
		S' \approx \int \left[L + \delta L\right] dt \approx \int \left[L + \dot{\varepsilon} \left( m \dot{\rv} \cdot \left(\rv - 2 \dot{\rv} t\right) + 2 t L\right)\right] \, dt
	\end{equation*}
	\begin{equation*}
		\delta L \approx \dot{\varepsilon}\left[ m \dot{\rv} \cdot \left(\rv - 2 \dot{\rv} t\right) + 2 t L\right]
	\end{equation*}
	Finalmente, la cantidad conservada será
	\begin{align*}
		Q &= \pd{\delta L}{\dot{\varepsilon}}\\
		&= m \dot{\rv} \cdot \left(\rv - 2 \dot{\rv} t\right) + 2 t L\\
		&= \pv \cdot \rv - 2 m \dot{\rv}^2 t + 2 t \left(\frac{1}{2}m \dot{\rv}^2 - U\right)\\
		&= \pv \cdot \rv - 2 t \left(\frac{1}{2}m \dot{\rv}^2 + U\right)\\
		&= \pv \cdot \rv - 2 H t
	\end{align*}
	
	Para comprobar que tanto $\n{M}$ como $Q$ se conservan, los momentos generalizados serán, usando \texttt{Mathematica}:
	\begin{equation*}
		p_r = m \dot{r}, ~~~~~~~ p_\theta = m r^2 \dot{\theta}, ~~~~~~~ p_\phi = m r^2 \dot{\phi} \sin^2 \theta
	\end{equation*}
	por lo que las ecuaciones de Euler-Lagrange:
	\begin{align*}
		m\ddot{r} &= m r \left(\dot{\phi}^2 \sin^2 \theta + \dot{\theta}^2 \right)+\frac{2 r U_0}{\left(r^2+\kappa  t\right)^2}\\
		\td{p_\theta}{t} &= m r^2 \dot{\phi}^2 \sin\theta \cos\theta\\
		\td{p_\phi}{t} &= \td{M_z}{t} = 0
	\end{align*}
	\qed
	
%%%%%%%%%%%%%%%%%%%%%%%%%%%%%%%%%%%%%%%%%%%%%%%%%%%%%%%%%%%%%%%%%%%%%%%%%%%%%%%%%%
	\item[(c)] We have seen in lectures that in static (time-independent) Coulomb potential the particle's trajectory always remains in the plane (so-called scattering plane). Analyze if this property remains valid in the time-dependent potential given in Eq. (\ref{eq:1.1}). Analyze if Kepler’s laws remain valid in this potential field.
	\resp El potencial es central y el momento angular $\n{M}$ es conservado componente por componente, por lo que el plano de rotación será constante.
	
	La segunda ley de Kepler se conservará porque el campo es central; la primera y tercera no lo harán, ya que dependen de una trayectoria específica que en este caso será alterada.
	
	\qed
	
\end{enumerate}

\section{Problema 2}
The action of the relativistic particle in the external spherically symmetric field is given by
\begin{equation}
	S = \int d\tau \, \Lambda, ~~~~~~~ \Lambda = \sqrt{g_{\mu\nu} \dot{x}^{\mu} \dot{x}^{\nu}}
\end{equation}
where $\tau$ is the proper time, ${x}^{\mu}(\tau)$ is the trajectory of the particle, $\dot{x}^{\mu} = \ltd{{x}^{\mu}}{\tau} \equiv u^\mu$ is the 4-velocity of the particle (we use $c = 1$ units as usual),
\begin{equation}
	g^{\mathrm{(sph)}}_{\mu\nu} = \mathrm{diag}\left\{a(r), ~ -b(r), ~ -r^2 b(r), ~ -r^2 b(r) \sin^2\theta\right\},
\end{equation}
\begin{equation}
	b(r) = \left(1 + \frac{A}{r}\right)^4, ~~~~~ a(r) = \frac{(r - A)^2}{(r + A)^2}, ~~~~~ A = \mathrm{const.} 
\end{equation}
and we assume that the particle moves only in the region $r > A > 0$.
\begin{enumerate}
%%%%%%%%%%%%%%%%%%%%%%%%%%%%%%%%%%%%%%%%%%%%%%%%%%%%%%%%%%%%%%%%%%%%%%%%%%%%%%%%%%
	\item[(a)] Identify all the continuous global symmetries of the action and construct the integrals of motion which may exist according to Noether's theorem.
	\resp Reemplazando explícitamente en $g_{\mu\nu}$
	\begin{equation*}
		g_{\mu\nu} = \begin{pmatrix}
			\frac{(r-A)^2}{(r + A)^2} & 0 & 0 & 0 \\
			0 & -\left(1 + \frac{A}{r}\right)^4 & 0 & 0 \\
			0 & 0 & - r^2 \left(1 + \frac{A}{r}\right)^4 & 0 \\
			0 & 0 & 0 & - r^2 \left(1 + \frac{A}{r}\right)^4 \sin^2 \theta
		\end{pmatrix}
	\end{equation*}
	Y desarrollando $\Lambda$
	\begin{align*}
		\Lambda &= \sqrt{\frac{(r-A)^2}{(r + A)^2} \dot{x}^{t}\dot{x}^{t} - \left(1 + \frac{A}{r}\right)^4 \dot{x}^{r}\dot{x}^{r} - \left(1 + \frac{A}{r}\right)^4 r^2 \dot{x}^{\theta}\dot{x}^{\theta} - \left(1 + \frac{A}{r}\right)^4 r^2 \sin^2 \theta \, \dot{x}^{\phi}\dot{x}^{\phi}}\\
		&= \sqrt{\frac{(r-A)^2}{(r + A)^2} \dot{x}^{t}\dot{x}^{t} - \left(1 + \frac{A}{r}\right)^4 \left[\dot{r}^2 + r^2 \dot{\theta}^2 + r^2 \sin^2 \theta \, \dot{\phi}^2 \right]}\\
		&= \sqrt{\frac{(r-A)^2}{(r + A)^2} \dot{x}^{t}\dot{x}^{t} - \left(1 + \frac{A}{r}\right)^4 \left[\dot{x}^2 + \dot{y}^2 + \dot{z}^2 \right]}
	\end{align*}
	Ya que la variable temporal no aparece explícitamente y el sistema tiene simetría esférica, tendremos conservación de la energía y el vector de momento angular $\n{M}$. Para la energía, usando \texttt{Mathematica}
	\begin{align*}
		\td{}{\tau} \pd{\Lambda}{\dot{t}} &= \pd{\Lambda}{t}\\
		\td{}{\tau} \left(\frac{a(r) \dot{t}}{\Lambda}\right) &= 0
	\end{align*}
	pero tenemos que se cumple que $\sqrt{g_{\mu\nu} \dot{x}^{\mu} \dot{x}^{\nu}} = 1$, y que $\dot{t} = \ltd{t}{\tau}$, por lo que
	\begin{equation*}
		\frac{(r - A)^2}{(r + A)^2} \td{t}{\tau} = \mathrm{const.}
	\end{equation*}
	Para la componente $z$ del momento angular tenemos el resultado típico,
	\begin{align}
		M_z = \pd{\Lambda}{\dot{\phi}} = r^2 \left(1 + \frac{A}{r}\right)^4 \dot{\phi} \sin^2\theta
	\end{align}
	\qed
	
%%%%%%%%%%%%%%%%%%%%%%%%%%%%%%%%%%%%%%%%%%%%%%%%%%%%%%%%%%%%%%%%%%%%%%%%%%%%%%%%%%
	\item[(b)] Analyze if Kepler's laws are valid for this system. Can we consider that the particle always moves in the plane in this field?
	\resp Ya que el potencial no es coulombiano, tendremos que no se cumplirán ni la primera ni la tercera ley de Kepler. En el caso de la segunda ley de Kepler, ya que es un problema de relatividad, la existencia del factor $\gamma$ hará imposible su cumplimiento.
	
	Ya que $\n{M}$ se conserva componente por componente, el movimiento de la partícula estará confinado a un plano.
	
	\qed
	
%%%%%%%%%%%%%%%%%%%%%%%%%%%%%%%%%%%%%%%%%%%%%%%%%%%%%%%%%%%%%%%%%%%%%%%%%%%%%%%%%%
	\item[(c)] Write out the equations of motion.
	\resp Si establecemos $\theta = \pi/2$, tenemos las integrales de movimiento
	\begin{equation*}
		E = \frac{(r - A)^2}{(r + A)^2} \td{t}{\tau}, ~~~~~ M =  r^2 \left(1 + \frac{A}{r}\right)^4 \dot{\phi}
	\end{equation*}
	\begin{equation*}
		\Rightarrow \td{t}{\tau} = E \frac{(r + A)^2}{(r - A)^2} , ~~~~~ \dot{\phi} =  \frac{M}{r^2 \left(1 + \frac{A}{r}\right)^4 }
	\end{equation*}
	Ya que sabemos que $\sqrt{g_{\mu\nu} \dot{x}^{\mu} \dot{x}^{\nu}} = 1$, tendremos
	\begin{equation}
		\dot{r}^2 = \frac{1}{\left(1 + \frac{A}{r}\right)^4} \left[E^2 \frac{{(r - A)^2}}{(r + A)^2} - \frac{M^2}{r^2 \left(1 + \frac{A}{r}\right)^4} - 1\right]
	\end{equation}
	\qed
	
%%%%%%%%%%%%%%%%%%%%%%%%%%%%%%%%%%%%%%%%%%%%%%%%%%%%%%%%%%%%%%%%%%%%%%%%%%%%%%%%%%
	\item[(d)] Integrate the equations of motion and find the trajectories, assuming that the particle always moves in a plane $\theta = \pi/2$ and at large distances from the center ($r \gg A$).
	\resp Aproximando a primer orden, usando \texttt{Mathematica}
	\begin{equation*}
		a(r) \approx 1 - \frac{4A}{r}, ~~~~~ b(r) \approx 1 + \frac{4A}{r}, ~~~~~ \frac{1}{a(r)} \approx 1 + \frac{4A}{r}, ~~~~~ \frac{1}{b(r)} \approx 1 - \frac{4A}{r}
	\end{equation*}
	En la ecuación radial del problema anterior,
	\begin{equation*}
		\dot{r}^2 = (E^2 - 1) + \frac{4 A}{r} - \frac{L^2}{r^2}
	\end{equation*}
	Solucionando con \texttt{Mathematica},
	\begin{equation*}
		\ldots
	\end{equation*}
\end{enumerate}








\end{document}





















